% Options for packages loaded elsewhere
\PassOptionsToPackage{unicode}{hyperref}
\PassOptionsToPackage{hyphens}{url}
%
\documentclass[
]{article}
\usepackage{amsmath,amssymb}
\usepackage{lmodern}
\usepackage{iftex}
\ifPDFTeX
  \usepackage[T1]{fontenc}
  \usepackage[utf8]{inputenc}
  \usepackage{textcomp} % provide euro and other symbols
\else % if luatex or xetex
  \usepackage{unicode-math}
  \defaultfontfeatures{Scale=MatchLowercase}
  \defaultfontfeatures[\rmfamily]{Ligatures=TeX,Scale=1}
\fi
% Use upquote if available, for straight quotes in verbatim environments
\IfFileExists{upquote.sty}{\usepackage{upquote}}{}
\IfFileExists{microtype.sty}{% use microtype if available
  \usepackage[]{microtype}
  \UseMicrotypeSet[protrusion]{basicmath} % disable protrusion for tt fonts
}{}
\makeatletter
\@ifundefined{KOMAClassName}{% if non-KOMA class
  \IfFileExists{parskip.sty}{%
    \usepackage{parskip}
  }{% else
    \setlength{\parindent}{0pt}
    \setlength{\parskip}{6pt plus 2pt minus 1pt}}
}{% if KOMA class
  \KOMAoptions{parskip=half}}
\makeatother
\usepackage{xcolor}
\IfFileExists{xurl.sty}{\usepackage{xurl}}{} % add URL line breaks if available
\IfFileExists{bookmark.sty}{\usepackage{bookmark}}{\usepackage{hyperref}}
\hypersetup{
  pdftitle={Comparison between t-SNE and UMAP},
  pdfauthor={Fan Cheng},
  hidelinks,
  pdfcreator={LaTeX via pandoc}}
\urlstyle{same} % disable monospaced font for URLs
\usepackage[margin=1in]{geometry}
\usepackage{longtable,booktabs,array}
\usepackage{calc} % for calculating minipage widths
% Correct order of tables after \paragraph or \subparagraph
\usepackage{etoolbox}
\makeatletter
\patchcmd\longtable{\par}{\if@noskipsec\mbox{}\fi\par}{}{}
\makeatother
% Allow footnotes in longtable head/foot
\IfFileExists{footnotehyper.sty}{\usepackage{footnotehyper}}{\usepackage{footnote}}
\makesavenoteenv{longtable}
\usepackage{graphicx}
\makeatletter
\def\maxwidth{\ifdim\Gin@nat@width>\linewidth\linewidth\else\Gin@nat@width\fi}
\def\maxheight{\ifdim\Gin@nat@height>\textheight\textheight\else\Gin@nat@height\fi}
\makeatother
% Scale images if necessary, so that they will not overflow the page
% margins by default, and it is still possible to overwrite the defaults
% using explicit options in \includegraphics[width, height, ...]{}
\setkeys{Gin}{width=\maxwidth,height=\maxheight,keepaspectratio}
% Set default figure placement to htbp
\makeatletter
\def\fps@figure{htbp}
\makeatother
\setlength{\emergencystretch}{3em} % prevent overfull lines
\providecommand{\tightlist}{%
  \setlength{\itemsep}{0pt}\setlength{\parskip}{0pt}}
\setcounter{secnumdepth}{5}
%% Any special functions or other packages can be loaded here.
\usepackage[utf8]{inputenc}
\usepackage[flushleft]{threeparttable}
\usepackage{multirow}
\usepackage[normalem]{ulem}
\useunder{\uline}{\ul}{}
%% \usepackage{makecell,interfaces-makecell} %latex error
\usepackage[]{graphicx}
\usepackage[]{color}
\usepackage{xcolor}
% \usepackage[ruled,vlined,linesnumbered]{algorithm2e}
\usepackage{algorithmic}
\usepackage{bm}
\usepackage{etoolbox}
\usepackage{tabu}
\usepackage{placeins}
\setlength {\marginparwidth }{2cm}
\usepackage{todonotes}
\urlstyle{tt}  % don't use monospace font for urls
\allowdisplaybreaks


\usepackage[ruled,vlined,linesnumbered]{algorithm2e}
\usepackage{amsmath, amsfonts, amssymb, bm, mathtools, todonotes, booktabs, microtype}

\makeatletter
%If bera loaded, this is monash wp.
\@ifpackageloaded{bera}{}{\usepackage{mathptmx}}
\makeatother
\usepackage[cal = txupr]{mathalpha}


\makeatletter
\@ifpackageloaded{geometry}{}{\usepackage{geometry}}
\makeatother
\geometry{margin=2.2cm}

\mathtoolsset{showonlyrefs}
\newcommand{\argmax}{\mathop{\text{arg\,max}}}
\clubpenalty = 10000
\widowpenalty = 10000
\brokenpenalty = 10000
\usepackage{graphicx}
\setcounter{topnumber}{2}
\setcounter{bottomnumber}{2}
\setcounter{totalnumber}{4}
\renewcommand{\topfraction}{0.9}
\renewcommand{\bottomfraction}{0.9}
\renewcommand{\textfraction}{0.1}
\renewcommand{\floatpagefraction}{0.9}

\usepackage{caption}
\DeclareCaptionStyle{italic}[justification=centering]
 {labelfont={bf},textfont={it},labelsep=colon}
\captionsetup[figure]{style=italic,format=hang,singlelinecheck=true}
\captionsetup[table]{style=italic,format=hang,singlelinecheck=true}

\ifLuaTeX
  \usepackage{selnolig}  % disable illegal ligatures
\fi

\title{Comparison between t-SNE and UMAP}
\author{Fan Cheng}
\date{2021-08-05}

\begin{document}
\maketitle

\hypertarget{t-sne}{%
\section{t-SNE}\label{t-sne}}

t-Distributed Stochastic Neighbor Embedding {[}t-SNE; @Van\_der\_Maaten2008-dv{]} is a dimension reduction technique well suited for the visualization of high-dimensional data in scatterplots. It works by minimizing the Kullback-Leibler divergence \(KL(P||Q)\) between two joint distributions, \(p_{ij}\) and \(q_{ij}\), measuring the pairwise similarities of the high-dimensional input data \(x_i\) and the low-dimensional embedded points \(y_i\) respectively.
To better capture the local data structure and solve the crowding problem, a normalized Student-t kernel is chosen for \(q_{ij}\) since it is heavy-tailed, while a Gaussian kernel is used for \(p_{ij}\). The bandwidth of the Gaussian kernel, \(\sigma_i\) is set to control the \emph{perplexity} of the conditional distribution, relating to the density of the input data. Considering the asymmetry of the Kullback-Leibler divergence, the object function aims at modelling similar points in the input space by nearby points in the embedding space. Analytical form of the gradient is used in the Gradient Descent optimization of Kullback-Leibler divergence. Details can be found in Algorithm \ref{alg:tsne}. Variants of the Barnes-Hut algorithm can be used to approximate the gradient of the object function {[}@Van\_Der\_Maaten2014-in{]}.

\begin{algorithm}[!htb]
  \caption{t-SNE}
  \label{alg:tsne}
  \DontPrintSemicolon
  \SetAlgoLined
  \SetKwInOut{Input}{Input}\SetKwInOut{Output}{Output}\SetKwInOut{Parameter}{parameter}
  \Input{ high-dimensional data $x_i$ for all $i=1,\ldots,N$ }
  \Output{ low-dimensional embedding $y_i$ for all $i=1,\ldots,N$ }
  \Parameter{ perplexity, number of iterations $T$, learning rate $\eta$, momentum $\alpha(t)$}
  \BlankLine
  \begin{algorithmic}[1]

  \STATE Compute the Gaussian probability of distances between $x_i$ and $x_j$ as $p_{j \mid i}$ using
  $$
    p_{j \mid i}=\frac{\exp \left(-\left\|x_{i}-x_{j}\right\|^{2} / 2 \sigma_{i}^{2}\right)}{\sum_{k \neq i} \exp \left(-\left\|x_{i}-x_{k}\right\|^{2} / 2 \sigma_{i}^{2}\right)};
  $$
  with the variance $\sigma_{i}$ is optimized by a constraint, $\text{perplexity} = 2 ^ {-\sum_j{p_{j \mid i} \log_2{p_{j \mid i}}}}$; 

  \STATE Set the probability as $p_{i j}=(p_{j \mid i}+p_{i \mid j}) / (2 N)$;

  \STATE Initialize the optimization solution as $\mathcal{Y}^{(0)} = \{ y_1, \dots, y_N \}$ sampled from normal distribution $\mathcal{N}(0, 10^{-4} \boldsymbol{I})$;

  \FOR{$t = 1,\dots,N$}

  \STATE Compute the Student t-distribution of distances between $y_i$ and $y_j$ as 
  $$
    q_{i j}=\frac{\left(1+\left\|y_{i}-y_{j}\right\|^{2}\right)^{-1}}{\sum_{k \neq l}\left(1+\left\|y_{k}-y_{l}\right\|^{2}\right)^{-1}};
  $$

  \STATE Compute the gradient of the Kullback-Leibler divergence loss function as
  $$
    \frac{\delta KL}{\delta y_{i}}=4 \sum_{j}\left(p_{i j}-q_{i j}\right)\left(y_{i}-y_{j}\right)\left(1+\left\|y_{i}-y_{j}\right\|^{2}\right)^{-1} ;
  $$

  \STATE Update $ \mathcal{Y}^{(t)}=\mathcal{Y}^{(t-1)}+\eta \frac{\delta KL}{\delta \gamma}+\alpha(t)\left(\mathcal{Y}^{(t-1)}-\mathcal{Y}^{(t-2)}\right). $

  \ENDFOR

  \end{algorithmic}
\end{algorithm}

Perplexity: to balance attention between local and global aspects of your data.

\begin{itemize}
\tightlist
\item
  The parameter is a guess about the number of close neighbors each point has. The original paper says, \emph{``The performance of SNE is fairly robust to changes in the perplexity, and typical values are between 5 and 50.''}
\end{itemize}

Other hyperparameters: these choices interact with the amount of momentum and the step size are employed in the gradient descent. It is therefore common to run the optimization several times on a data set to find appropriate values for the parameters.

\hypertarget{illustration}{%
\subsection{Illustration}\label{illustration}}

Below are some figures to illustrate the steps in t-SNE from a book named \emph{Machine Learning with R, the tidyverse, and mlr} by Hefin I. Rhys.

\begin{figure}

{\centering \includegraphics[width=0.8\linewidth]{/Users/fche0019/git/kderm/reports/figures/tsne1} 

}

\end{figure}
\begin{figure}

{\centering \includegraphics[width=0.8\linewidth]{/Users/fche0019/git/kderm/reports/figures/tsne2} 

}

\end{figure}
\begin{figure}

{\centering \includegraphics[width=0.8\linewidth]{/Users/fche0019/git/kderm/reports/figures/tsne3} 

}

\end{figure}
\begin{figure}

{\centering \includegraphics[width=0.8\linewidth]{/Users/fche0019/git/kderm/reports/figures/tsne4} 

}

\end{figure}

\hypertarget{umap}{%
\section{UMAP}\label{umap}}

UMAP {[}Uniform Manifold Approximation and Projection; @McInnes2018-xo{]} is a competitive dimension reduction methods with t-SNE in increasing speed and preseving global structure. It is based on Laplacian Eigenmaps, considering Riemannian geometry and algebraic topology. Instead of optimizing the similarities in t-SNE, UMAP works by minimizing the error between two topological representations from the input space and embedding space. The error is measured by the cross entropy \(C\) of two local fuzzy simplical sets built from the nearest neighborhood graph, with edge weights representing the likelihood that two points are connected.
UMAP extends a raidus outwards at each point until overlap with the nearest neighbors, with a \(min_dist\) as a desired seperation between nearby embedded points to control global structure. Finally, similar to t-SNE, stochastic gradient descent is used for the optimization of object function. The algorithm and hyperparameters are detailed in Section 4 of @McInnes2018-xo.

\hypertarget{parameters}{%
\subsection{Parameters}\label{parameters}}

\begin{itemize}
\item
  K: number of nearest neighbors
\item
  min\_dist: the minimum distance between points in low-dimensional space.

  This parameter controls how tightly UMAP clumps points together, with low values leading to more tightly packed embeddings. Larger values of min\_dist will make UMAP pack points together more loosely, focusing instead on the preservation of the broad topological structure.
\end{itemize}

\hypertarget{compare}{%
\section{Compare}\label{compare}}

The biggest difference between the the output of UMAP when compared with t-SNE is this balance between local and global structure - UMAP is often better at preserving global structure.

\hypertarget{visualization-examples}{%
\subsection{Visualization examples}\label{visualization-examples}}

\url{https://pair-code.github.io/understanding-umap/}

\hypertarget{t-sne-disadvantages}{%
\subsection{t-SNE disadvantages}\label{t-sne-disadvantages}}

\begin{itemize}
\item
  tSNE does not scale well for rapidly increasing sample sizes.
\item
  tSNE does not preserve global data structure. Only within cluster distances are meaningful while between cluster similarities are not guaranteed, therefore it is widely acknowledged that clustering on tSNE is not a very good idea.
\item
  tSNE can practically only embed into 2 or 3 dimensions.
\item
  tSNE can not work with high-dimensional data directly, Autoencoder or PCA are often used for performing a pre-dimensionality reduction before plugging it into the tSNE.
\item
  tSNE consumes too much memory for its computations which becomes especially obvious when using large perplexity.
\end{itemize}

\hypertarget{takeaways-from-the-blog}{%
\subsection{Takeaways from the blog}\label{takeaways-from-the-blog}}

\begin{enumerate}
\def\labelenumi{\arabic{enumi}.}
\item
  Hyperparameters really matter
\item
  Cluster sizes in a UMAP plot mean nothing
\item
  Distances between clusters might not mean anything
\item
  Random noise doesn't always look random
\end{enumerate}

Especially at low values of n\_neighbors, spurious clustering can be observed.

\begin{enumerate}
\def\labelenumi{\arabic{enumi}.}
\setcounter{enumi}{4}
\tightlist
\item
  You may need more than one plot
\end{enumerate}

\hypertarget{things-related-to-our-paper}{%
\section{Things related to our paper}\label{things-related-to-our-paper}}

t-SNE Variants: implemented via Barnes-Hut approximations for gradient descent process, which also limits the implementation to two or three dimensions. @Tang2016-ho extended t-SNE with the LargeVis algorithm where random projection trees are used to reduce time cost.

L.J.P. van der Maaten. Accelerating t-SNE using Tree-Based Algorithms. Journal of Machine Learning Research 15(Oct):3221-3245, 2014.

UMAP: allow any distance metric for constructing high-dimensional distance probabilities using NN-Descent {[}@Dong2011-ee{]} algorithm; used procrustes measure to measure stability of the embedding.

\hypertarget{estimating-intrinsic-dimensions}{%
\section{Estimating intrinsic dimensions}\label{estimating-intrinsic-dimensions}}

There are multiple methods to estimate the intrinsic dimensions, with 17 of them implemented in R package \emph{Rdimtools}.

\url{https://kisungyou.com/Rdimtools/reference/index.html\#section--dimension-estimation}

\begin{itemize}
\item
  Running 5 of them on MNIST dataset gives around 12 dimensions as an output.
\item
  How to deal with the implementation of t-SNE then?
\item
  How to unify this in the paper?
\end{itemize}

\end{document}
